برای حل این مسئله، کاردینال مجموعه اعداد حقیقی را با $c$ یا $2^{\aleph_0}$ نمایش می‌دهیم، که در آن $\aleph_0$ کاردینال مجموعه اعداد طبیعی است.


\subsectionaddtolist{الف}

برای اثبات این حکم، ابتدا نشان می‌دهیم که $|\mathbb{R}| = |\mathbb{R}^2|$ و سپس با استفاده از استقرا آن را به هر $n$ طبیعی تعمیم می‌دهیم.

می‌دانیم که کاردینال اعداد حقیقی برابر است با:
\[ |\mathbb{R}| = 2^{\aleph_0} \]

حال کاردینال $\mathbb{R}^2$ را محاسبه می‌کنیم:
\[ |\mathbb{R}^2| = |\mathbb{R} \times \mathbb{R}| = |\mathbb{R}| \cdot |\mathbb{R}| = 2^{\aleph_0} \cdot 2^{\aleph_0} = 2^{\aleph_0 + \aleph_0} \]

از آنجایی که اجتماع دو مجموعه شمارا، باز هم شمارا است ($\aleph_0 + \aleph_0 = \aleph_0$)، پس:
\[ 2^{\aleph_0 + \aleph_0} = 2^{\aleph_0} = |\mathbb{R}| \]

بنابراین:
\[ |\mathbb{R}^2| = |\mathbb{R}| \]

\begin{itemize}
    \item {پایه استقرا ($n=1$):}
    \[ |\mathbb{R}| = |\mathbb{R}^1| \]
    
    \item {فرض استقرا:} فرض می‌کنیم برای $n=k$ حکم برقرار باشد، یعنی:
    \[ |\mathbb{R}^k| = |\mathbb{R}| \]
    
    \item {گام استقرا ($n=k+1$):} باید نشان دهیم $|\mathbb{R}^{k+1}| = |\mathbb{R}|$. طبق تعریف حاصل‌ضرب دکارتی داریم:
    \[ |\mathbb{R}^{k+1}| = |\mathbb{R}^k \times \mathbb{R}| = |\mathbb{R}^k| \cdot |\mathbb{R}| \]
    
    با استفاده از فرض استقرا، جایگذاری می‌کنیم:
    \[ |\mathbb{R}^k| \cdot |\mathbb{R}| = |\mathbb{R}| \cdot |\mathbb{R}| = |\mathbb{R}^2| \]
    
    طبق گام اول می‌دانیم $|\mathbb{R}^2| = |\mathbb{R}|$، پس:
    \[ |\mathbb{R}^{k+1}| = |\mathbb{R}| \]
\end{itemize}
بنابراین حکم برای هر عدد طبیعی $n$ برقرار است.

\newpage

\subsectionaddtolist{ب}

فرض کنید $\{A_i\}_{i=1}^{\infty}$ مجموعه‌ای شمارا از مجموعه‌ها باشد به طوری که برای هر $i$ داشته باشیم $|A_i| = |\mathbb{R}| = c$. می‌خواهیم ثابت کنیم:
\[ \left| \bigcup_{i=1}^{\infty} A_i \right| = c \]


 از آنجایی که هر $A_i$ زیرمجموعه‌ای از این اجتماع است (مثلاً $A_1 \subseteq \bigcup_{i=1}^{\infty} A_i$)، کاردینال اجتماع نمی‌تواند کمتر از کاردینال یکی از مجموعه‌های درون آن باشد:
    \[ \left| \bigcup_{i=1}^{\infty} A_i \right| \ge |A_1| = c \]
    
 این اجتماع را می‌توان به صورت نگاشت فرارونده‌ای از حاصل‌ضرب دکارتی اندیس‌ها و مجموعه‌ها، یعنی زیرمجموعه‌ای از $\mathbb{N} \times \mathbb{R}$ در نظر گرفت. بنابراین:
    \[ \left| \bigcup_{i=1}^{\infty} A_i \right| \le |\mathbb{N} \times \mathbb{R}| = |\mathbb{N}| \cdot |\mathbb{R}| = \aleph_0 \cdot c \]
    
    از قوانین حساب کاردینال‌ها می‌دانیم حاصل‌ضرب یک کاردینال نامتناهی در کاردینالی بزرگ‌تر یا مساوی خودش، برابر با کاردینال بزرگ‌تر است ($\aleph_0 \cdot c = c$). پس:
    \[ \left| \bigcup_{i=1}^{\infty} A_i \right| \le c \]

با توجه به اینکه کاردینال این اجتماع هم بزرگ‌تر یا مساوی $c$ و هم کوچک‌تر یا مساوی $c$ شد، نتیجه می‌گیریم:
\[ \left| \bigcup_{i=1}^{\infty} A_i \right| = c \]

\subsectionaddtolist{ج}

مجموعه تمام توابع از $\mathbb{R}$ به $\mathbb{R}$ را با $\mathbb{R}^\mathbb{R}$ نمایش می‌دهند. می‌خواهیم نشان دهیم $|\mathbb{R}^\mathbb{R}| > |\mathbb{R}|$.

ابتدا کاردینال این مجموعه توابع را محاسبه می‌کنیم:
\[ |\mathbb{R}^\mathbb{R}| = |\mathbb{R}|^{|\mathbb{R}|} = (2^{\aleph_0})^{c} \]

با استفاده از قوانین توان‌ها داریم:
\[ (2^{\aleph_0})^{c} = 2^{\aleph_0 \cdot c} \]

از آنجایی که $c$ یک کاردینال نامتناهی و بزرگ‌تر از $\aleph_0$ است، حاصل‌ضرب $\aleph_0 \cdot c$ برابر با خود $c$ می‌شود:
\[ 2^{\aleph_0 \cdot c} = 2^c \]

بنابراین، کاردینال مجموعه توابع برابر با $2^c$ (کاردینال مجموعه توانی $\mathbb{R}$) است. برای هر مجموعه $X$، کاردینال مجموعه توانی آن همیشه اکیداً بزرگ‌تر از کاردینال خود مجموعه است:
\[ 2^{|X|} > |X| \]

با قرار دادن $X = \mathbb{R}$ خواهیم داشت:
\[ 2^c > c \]

که نتیجه می‌دهد:
\[ |\mathbb{R}^\mathbb{R}| > |\mathbb{R}| \]
